\documentclass[UTF8]{ctexart}

% 支持从项目根目录、概率论与数理统计目录或当前文档目录读取共享样式。
\makeatletter
\def\input@path{{../}{../../}}
\makeatother
\usepackage{researchnotes}

\title{数学期望性质入门}
\author{}
\date{}

\begin{document}
\maketitle

\section{从平均数认识数学期望}
\label{sec:basic-expectation-introduction}

一种彩票有时中奖、有时不中，只看最高奖金，不能判断平均能得到多少。
既要考虑每种结果的数值，也要考虑它出现的概率。
\keyterm{数学期望}（mathematical expectation）就是\keyterm{以概率为权重的平均数}。

本文只讨论可能取值有限的情况，因此所有公式中的期望都是有限数，
只需概率、加权平均和初等代数知识。
先理解期望的含义，再学习常用运算规则，最后用“分组平均”认识条件期望和全期望公式。

\subsection{随机变量与期望的计算}

\keyterm{随机变量}（random variable）是给随机结果对应一个数。
例如，掷骰子后朝上的点数可以记为 $X$。
大写 $X$ 表示结果尚未确定时的随机量，小写 $x$ 表示一次具体取到的数值。
记号 $\mathbb P(X=x)$ 表示“$X$ 取到 $x$ 的概率”。

\begin{definition}[有限个取值时的数学期望]
设 $X$ 的不同可能取值为 $x_1,\ldots,x_n$，相应概率为 $p_1,\ldots,p_n$，
其中 $p_i\geq0$ 且 $p_1+\cdots+p_n=1$。定义
\begin{equation}
  \mathbb E[X]=x_1p_1+\cdots+x_np_n
  =\sum_{i=1}^{n}x_ip_i.
  \label{eq:basic-expectation-definition}
\end{equation}
记号 $\mathbb E[X]$ 读作“$X$ 的期望”，$\sum$ 表示把列出的各项相加。
\end{definition}

计算时只做三步：列出可能取值，写出对应概率，再将“取值乘概率”的结果相加。
只有各种结果等可能时，才可以直接把这些取值相加后除以取值个数。

\begin{example}[公平骰子的点数]
公平六面骰子的点数 $1,2,\ldots,6$ 各以 $1/6$ 的概率出现，因此
\[
  \mathbb E[X]
  =1\times\frac16+\cdots+6\times\frac16
  =\frac{21}{6}=3.5.
\]
骰子不可能掷出 $3.5$ 点，这并不矛盾：
\keyterm{期望是平均水平，不是下一次必然发生的结果，也不必是某个实际取值。}
\end{example}

\subsection{期望与多次试验的平均值}

若独立地重复同一种试验，且每次各种结果的概率保持不变，
试验次数足够多时，实际平均值以很高的概率接近期望。
这一联系由\keyterm{大数定律}（law of large numbers）保证。
它不表示少量试验的平均值一定接近期望，也不表示每次增加一次试验，误差都必然变小。

\begin{example}[奖金与净收益]
\label{ex:basic-expectation-lottery}
假设一次抽奖有 $70\%$ 的概率得 $0$ 元、$20\%$ 的概率得 $10$ 元、
$10\%$ 的概率得 $50$ 元。记奖金为 $X$，则
\[
  \mathbb E[X]=0\times0.7+10\times0.2+50\times0.1=7\text{ 元}.
\]
若每次参与需支付 $5$ 元，净收益是奖金减去费用，其期望为 $7-5=2$ 元。
这说明平均净收益为正，并不保证某一次不会亏钱；
例如未中奖时，该次净收益就是 $-5$ 元。
\end{example}

\section{最常用的运算规则：线性性}
\label{sec:basic-expectation-linearity}

\subsection{常数、倍数与固定增减}

\begin{property}[常数与一次式]
若 $a,b,c$ 是固定常数，则
\begin{equation}
  \mathbb E[c]=c,\qquad
  \mathbb E[aX+b]=a\mathbb E[X]+b.
  \label{eq:basic-expectation-affine}
\end{equation}
\end{property}

\begin{proof}
常数 $c$ 以概率 $1$ 出现，所以 $\mathbb E[c]=c$。
若 $X$ 取 $x_i$ 的概率为 $p_i$，则
\[
  \mathbb E[aX+b]
  =\sum_i(ax_i+b)p_i
  =a\sum_i x_ip_i+b\sum_i p_i
  =a\mathbb E[X]+b.
\]
即使变换后某些取值相同，把相同结果的概率合并也不会改变这份加权和。
\end{proof}

它的直观含义是：每次得分都翻倍，平均得分也翻倍；每次固定多得 $3$ 分，平均分也多 $3$ 分。
例 \ref{ex:basic-expectation-lottery} 中“净收益的期望等于奖金期望减去费用”，正是这一性质。
其中 $a$ 可以为负数，例如 $\mathbb E[-X]=-\mathbb E[X]$。

\subsection{和与差的期望}

\begin{property}[期望的线性性]
\label{prop:basic-expectation-linearity}
对随机变量 $X,Y$ 和常数 $a,b$，有
\begin{equation}
  \mathbb E[aX+bY]=a\mathbb E[X]+b\mathbb E[Y].
  \label{eq:basic-expectation-linearity}
\end{equation}
特别地，
$\mathbb E[X+Y]=\mathbb E[X]+\mathbb E[Y]$，
$\mathbb E[X-Y]=\mathbb E[X]-\mathbb E[Y]$。
这些公式\keyterm{不要求 $X,Y$ 独立}。
\end{property}

\begin{proof}
把同一次试验中 $(X,Y)$ 的可能结果依次列出：
第 $i$ 种结果为 $(x_i,y_i)$，概率为 $q_i$。
按所有配对结果求加权平均，有
\[
  \mathbb E[aX+bY]
  =\sum_i(ax_i+by_i)q_i
  =a\sum_i x_iq_i+b\sum_i y_iq_i.
\]
后面两个加权和分别为 $\mathbb E[X]$ 和 $\mathbb E[Y]$；
同一个 $X$ 取值即使出现在多个配对结果中，合并概率后仍是原来的期望。
因此无需假设独立。
\end{proof}

例如，掷出一次骰子的点数为 $X$，额外奖励为 $Y=2X$。
$Y$ 完全由 $X$ 决定，但仍有
$\mathbb E[X+Y]=3.5+7=10.5$。
可见，“不独立”不会破坏加法公式。

对有限个随机变量同样有
\[
  \mathbb E[X_1+\cdots+X_n]
  =\mathbb E[X_1]+\cdots+\mathbb E[X_n].
\]
实际求总分、总费用、总次数时，往往不必先求出总量的全部概率分布，
可以先求各部分的期望，再相加。

\section{大小关系、乘积与平方}
\label{sec:basic-expectation-other-properties}

\subsection{平均不会超出原来的范围}

\begin{property}[非负性与大小关系]
如果对每种可能结果都有 $X\geq0$，则 $\mathbb E[X]\geq0$。
如果每种可能结果下都有 $X\leq Y$，则 $\mathbb E[X]\leq\mathbb E[Y]$。
特别地，如果 $m\leq X\leq M$，其中 $m,M$ 为常数，则
\[
  m\leq\mathbb E[X]\leq M.
\]
\end{property}

\begin{proof}
非负数乘以非负概率，再相加，结果仍然非负。
若 $X\leq Y$，则 $Y-X\geq0$，所以
$\mathbb E[Y]-\mathbb E[X]=\mathbb E[Y-X]\geq0$。
最后分别比较 $m$ 与 $X$、$X$ 与 $M$ 即可。
\end{proof}

例如，成绩始终在 $0$ 到 $100$ 分之间，期望成绩也在这个范围内。
如果某种加分规则对每个学生都不减分，那么实施后的期望成绩不会下降。
这要求比较同一种结果下的两个数，不能只比较它们的最大值。

另一个有用的结果是
\[
  |\mathbb E[X]|\leq\mathbb E[|X|].
\]
这是因为 $-|X|\leq X\leq|X|$。
先平均时正负数可能相互抵消；先取绝对值再平均，就没有这种抵消。

\subsection{乘积需要额外条件}

有限取值情形下，称 $X,Y$ \keyterm{相互独立}（independent），
是指对所有可能的 $x,y$，都有
$\mathbb P(X=x,Y=y)=\mathbb P(X=x)\mathbb P(Y=y)$。
直观地说，知道其中一个的结果，不会改变另一个各种取值的概率。

\begin{property}[独立变量乘积的期望]
\label{prop:basic-expectation-product}
如果 $X,Y$ 独立，则
\begin{equation}
  \mathbb E[XY]=\mathbb E[X]\mathbb E[Y].
  \label{eq:basic-expectation-product}
\end{equation}
\end{property}

这个公式来自独立性：一对结果 $(x,y)$ 的概率可以拆成两个概率相乘。
因此，对所有配对结果计算“$xy$ 乘以对应概率”时，加权和可以因式分解为
\[
  \sum_x\sum_y xy\,\mathbb P(X=x)\mathbb P(Y=y)
  =\left(\sum_x x\mathbb P(X=x)\right)
   \left(\sum_y y\mathbb P(Y=y)\right).
\]
两个求和符号只是表示遍历所有配对结果。
例如，独立掷两枚公平骰子，点数分别为 $X,Y$，
则乘积的期望为 $3.5\times3.5=12.25$。

\begin{example}[不独立时不能直接拆乘积]
\label{ex:basic-expectation-product-failure}
抛一枚公平硬币，正面记 $X=1$，反面记 $X=-1$，并令 $Y=X$。
两者总是相等，因而不独立。此时
\[
  \mathbb E[X]=\mathbb E[Y]=0,\qquad
  \mathbb E[XY]=\mathbb E[X^2]=1.
\]
所以 $\mathbb E[XY]\neq\mathbb E[X]\mathbb E[Y]$。
独立性保证乘积公式成立；若不知道是否独立，就不能直接套用。
\end{example}

\subsection{先平方还是先求平均}

例 \ref{ex:basic-expectation-product-failure} 同时说明，
$\mathbb E[X^2]$ 和 $(\mathbb E[X])^2$ 一般不相等。
它们的运算顺序不同：前者先平方再平均，后者先平均再平方。

记 $\mu=\mathbb E[X]$。随机变量的\keyterm{方差}（variance）定义为
\begin{equation}
  \operatorname{Var}(X)
  =\mathbb E[(X-\mu)^2]
  =\mathbb E[X^2]-\mu^2.
  \label{eq:basic-expectation-variance}
\end{equation}
第一式表示“与均值的距离平方的平均”，用于描述结果的波动程度。
第二式由展开平方得到：
$\mathbb E[X^2-2\mu X+\mu^2]
=\mathbb E[X^2]-2\mu^2+\mu^2$。
由于平方非负，
\[
  \mathbb E[X^2]\geq(\mathbb E[X])^2.
\]
对于上述硬币例子，均值为 $0$，方差为 $1$：
平均得分为零，并不表示每次的得分都为零。

\section{把概率变成期望：零一变量}
\label{sec:basic-expectation-indicators}

如果某个事件发生时记为 $1$，不发生时记为 $0$，就得到一个
\keyterm{指示随机变量}（indicator random variable）$I$。
若事件发生的概率为 $p$，则
\begin{equation}
  \mathbb E[I]=1\times p+0\times(1-p)=p.
  \label{eq:basic-expectation-indicator}
\end{equation}
所以，\keyterm{零一变量的期望就是取到 $1$ 的概率}。

\begin{example}[猜选择题的期望答对数]
假设有 $10$ 道选择题，每题猜对的概率都是 $1/4$。
令 $I_i$ 表示第 $i$ 题是否答对，答对为 $1$，答错为 $0$。
答对总数为 $N=I_1+\cdots+I_{10}$，因此
\[
  \mathbb E[N]
  =\sum_{i=1}^{10}\mathbb E[I_i]
  =10\times\frac14=2.5.
\]
这里不要求不同题的答对情况独立，因为使用的是加法规则。
一次考试的答对数一定是整数；$2.5$ 描述的是期望答对数。
\end{example}

\section{条件期望与分组平均}
\label{sec:basic-expectation-conditioning}

\subsection{知道部分信息后，重新计算平均值}

\keyterm{条件期望}（conditional expectation）是在已经知道某个条件的情况下，
使用该条件下的概率重新求平均。
例如，$\mathbb E[X\mid Y=y]$ 表示“已知 $Y=y$ 时，$X$ 的期望”，
竖线 $\mid$ 可以读作“在已知……的条件下”。
本文只对发生概率为正的条件进行这种计算。

\begin{example}[先选抽奖箱，再抽奖金]
\label{ex:basic-expectation-total}
假设先随机选择一个抽奖箱：以 $40\%$ 的概率选 A 箱，以 $60\%$ 的概率选 B 箱。
A 箱内有两张分别标着 $0$ 元、$10$ 元的奖券，B 箱内有两张分别标着 $10$ 元、$20$ 元的奖券；
选好箱子后，在其中等概率抽取一张。记奖金为 $X$，选中的箱子为 $Y$。

已知选中哪个箱子时，分别有
\[
  \mathbb E[X\mid Y=\mathrm A]=\frac{0+10}{2}=5,\qquad
  \mathbb E[X\mid Y=\mathrm B]=\frac{10+20}{2}=15.
\]
尚未知道箱子时，按选箱概率加权，得到
\[
  \mathbb E[X]=0.4\times5+0.6\times15=11\text{ 元}.
\]
也可以把所有抽取路径直接展开验证：
来自 A 箱的两张券各以 $0.2$ 的概率被抽到，
来自 B 箱的两张券各以 $0.3$ 的概率被抽到，因此
$0\times0.2+10\times0.2+10\times0.3+20\times0.3=11$。
\end{example}

\subsection{全期望公式：先在组内平均，再在组间平均}

设分组变量 $Y$ 的取值为 $y_1,\ldots,y_m$，每组概率为
$q_j=\mathbb P(Y=y_j)>0$，则
\begin{equation}
  \mathbb E[X]
  =\sum_{j=1}^{m}
    \mathbb E[X\mid Y=y_j]q_j.
  \label{eq:basic-expectation-total}
\end{equation}
这称为\keyterm{全期望公式}（law of total expectation）。
每个可能结果只属于一个组；先在各组内求平均，再按组的概率加权，
每个结果最终获得的权重仍是它原来的发生概率，所以与直接求平均一致。

例 \ref{ex:basic-expectation-total} 中不能直接计算 $(5+15)/2=10$，
因为两个箱子被选中的概率并不相等。
类似地，把各班平均分合成全年级平均分时，应按各班人数占比加权；
只有各班人数相等时，才可直接对班级平均分作算术平均。

\subsection{为什么还能写成两层期望}

记 $M=\mathbb E[X\mid Y]$，它表示“根据实际选中的组，返回该组的平均值”。
在抽奖箱例子中，$M$ 以 $0.4$ 的概率等于 $5$、以 $0.6$ 的概率等于 $15$，
所以 $M$ 本身仍然是随机变量。
再求一次期望，有
\begin{equation}
  \mathbb E[\mathbb E[X\mid Y]]
  =\mathbb E[M]
  =\mathbb E[X].
  \label{eq:basic-expectation-tower}
\end{equation}
内层给出各组的平均值，外层把这些平均值按组的概率加权。

这不表示 $X=M$。在 A 箱中，实际奖金 $X$ 是 $0$ 或 $10$，
而该组均值 $M$ 是 $5$，两者每次都不相等；
但是组内和总体上的平均值都相同。

进一步，若已经知道 $Y$，又按额外信息 $Z$ 细分，就有
\[
  \mathbb E[\mathbb E[X\mid Y,Z]\mid Y]
  =\mathbb E[X\mid Y].
\]
这称为\keyterm{迭代期望法则}（law of iterated expectations），也称塔式法则。
例如，从全校学生中等概率抽取一人，令 $X$ 为其成绩，$Y$ 为其年级，$Z$ 为其班级：
在固定年级内，先算各班平均分，再按班级人数加权，
与直接计算这个年级的平均分相同。
内层保留了外层已经知道的年级信息，不能随意改成只根据一个无关条件求平均。

\section{常用性质速查与做题顺序}
\label{sec:basic-expectation-summary}

表 \ref{tab:basic-expectation-summary} 汇总本文最常用的公式。
它们都在本文“随机变量只有有限个可能取值”的范围内使用。

\begin{table}
  \centering
  \caption{数学期望常用性质}
  \label{tab:basic-expectation-summary}
  \begin{tabularx}{\textwidth}{lYY}
    \toprule
    性质 & 公式 & 使用提醒 \\
    \midrule
    常数 & $\mathbb E[c]=c$ & $c$ 是固定的数 \\
    固定倍数与增减 & $\mathbb E[aX+b]=a\mathbb E[X]+b$ & $a,b$ 都是常数 \\
    加法 & $\mathbb E[X+Y]=\mathbb E[X]+\mathbb E[Y]$ & 不要求独立 \\
    大小关系 & $X\leq Y\Rightarrow\mathbb E[X]\leq\mathbb E[Y]$ & 每种可能结果下都要满足大小关系 \\
    乘积 & $\mathbb E[XY]=\mathbb E[X]\mathbb E[Y]$ & 独立时可用，一般不能直接拆分 \\
    零一变量 & $\mathbb E[I]=\mathbb P(I=1)$ & $I$ 只能取 $0,1$ \\
    平方 & $\mathbb E[X^2]\geq(\mathbb E[X])^2$ & 两种运算顺序一般不相等 \\
    分组平均 & $\mathbb E[\mathbb E[X\mid Y]]=\mathbb E[X]$ & 各组均值按组的概率加权 \\
    \bottomrule
  \end{tabularx}
\end{table}

做题时，可以依次考虑：
\begin{enumerate}
  \item \textbf{先确定要求什么。} 是奖金、净收益、答对数，还是若干步的总得分？
  \item \textbf{取值少时，直接列分布。} 检查概率之和是否为 $1$，再按定义计算。
  \item \textbf{总量能拆成和时，优先用线性性。} 各部分即使相关，期望仍能相加。
  \item \textbf{遇到乘积时，先检查独立性。} 遇到平方、商等运算，不直接把期望移进去。
  \item \textbf{已知某个条件后更容易计算时，尝试分组。} 先求各组条件期望，再按正确概率加权。
  \item \textbf{最后检查含义与范围。} 期望是否落在可能值的最小值与最大值之间？
  是否误把平均收益解释成了每次都能得到的收益？
\end{enumerate}

\end{document}
